% This reponse letter was taken from Vilens' paper
% The original files are by David Wisth

\documentclass[a4paper,11pt]{article}
\usepackage[utf8]{inputenc}
\usepackage{csquotes}
\usepackage[margin=3cm]{geometry}
\pagenumbering{gobble} % remove page numbering
\usepackage[dvipsnames]{xcolor}
\renewcommand{\mkbegdispquote}[2]{\color{Blue}}

%\usepackage{todonotes}
\usepackage[final]{pdfpages}
\setlength\parindent{0pt} % remove indent on paragraphs

\usepackage{amsmath} % assumes amsmath package installed
\usepackage{amssymb}  % assumes amsmath package installed
\usepackage{siunitx} 

\usepackage{outlines}

\input{_math_notation.tex}

\usepackage[color=yellow]{todonotes}
\presetkeys{todonotes}{inline}{}
%\newcommand\todo[1]{\textbf{[TODO: #1}]}

%opening
\title{Authors' Response to Editorial Staff and Reviewers}
\author{}
\date{}

\begin{document}

\maketitle
\vspace{-5ex}

\begin{displayquote}
Dear Mr. Mat\'ias Mattamala,

Your paper

21-2269

"An Efficient Locally Reactive Controller for Safe Navigation in Visual
Teach and Repeat Missions" by Mat\'ias Mattamala, Nived Chebrolu, Maurice Fallon

Submission for RA-L and ICRA submitted to the IEEE Robotics and Automation 
Letters (RA-L) has been reviewed by the Associate Editor and selected 
Reviewers. The reviews of the paper are attached.

On the basis of the reviewers' ratings and comments, we decided that your paper 
cannot be published in the Letters in the present form. 
However, you are encouraged to rewrite and resubmit a revised version
of your work addressing all editorial concerns.

[Deleted instructions related to resubmission]

If you have any related concern, do not hesitate to contact me.

Sincerely,

Eric Marchand\\
Editor (Vision and Sensor-Based Control)\\
IEEE Robotics and Automation Letters
\end{displayquote}

Dear Prof. Marchand,\\
\\
We would like to thank all the editors and reviewers for the time taken 
to review this paper and providing many valuable comments and suggestions.
We have addressed these comments in our resubmission and a 
document showing the difference between the original and revised versions of 
the paper is appended to this response.\\
\\
Yours sincerely,\\
\\
Mat\'ias Mattamala (on behalf of the authors)

%%%%%%%%%%%%%%%%%%%%%%%%%%%%%%%%%%%%%%%%%%%%%%%%%%%%%%%%%%%%%%%%%%%%%%%%%%%%%%%%
%%%%%%%%%%%%%%%%%%%%%%%%%%%%%%%%%%%%%%%%%%%%%%%%%%%%%%%%%%%%%%%%%%%%%%%%%%%%%%%%

\pagebreak
\section*{Response to Editor}

\begin{displayquote}
  All the Reviewers agree that the paper is interesting and very relevant
to the field. Nonetheless, some concerns have been raised, in
particular: 

- The experimental section does not provide any comparison with
state-of-the-art works. The authors are encouraged to compare their
approach against other baselines, or at least provide a justification
about why their work is theoretically superior to others;
\end{displayquote}

First of all, thanks for the excellent summary of the main concerns. 
\begin{outline}[enumerate]
\1 Regarding  the baselines, we have added comparisons to three local 
planner/controller baselines:

	\2 Sampling-based local planner adapted from open code by Zhang et 
	al. [37].
	\2 Potential field controller, which is a special case of our proposed 
	controller (\emph{Obstacle-free Goal Reaching RMP} and \emph{Obstacle 
	Avoidance RMP} with fields always enabled).
	\2 Geodesic Distance Field (GDF) based controller, which is another 
	special case (\emph{Obstacle-free Goal Reaching RMP} and \emph{GDF-based 
	Goal Reaching RMP} with both fields always enabled).

\1 We compared the methods in simulation recording a reference path as done 
in a Visual Teach and Repeat scheme. Then we introduced some artificial 
obstacles close to the recorded path, similarly to the Experiment 1 reported in 
the paper in the indoor environment. 
\1 Then we ran each controller 10 times using a 
moving look-ahead goal as input to record the tracked trajectories and number 
of collisions.
\1 Our results show that our approach allows the robot to track the trajectory 
staying away from obstacles, while other methods perform poorly, frequently 
hitting obstacles or getting stuck in local minima.
\1 The experiment was also added to the supplementary video to illustrate the 
previous points.
\end{outline}

%\begin{outline}[enumerate]
%	\item Added comparisons to three baseline approaches 
%	(a) Sampling based approach (add cite)
%	(b) Potential field based controller (add cite),
%	(c) Geodesic distance field (GDF) based controller (add cite)
%	\item Provide a short summary of the baseline experiment results (1-2 
%lines).
%\end{outline}

\begin{displayquote}
- There are some concepts and design choices that should be better
explained and motivated;
\end{displayquote}

We have made some clarifications to address the concerns about the parameters, 
basic RMP concepts, and motivation for the framework. In particular:
\begin{outline}[enumerate]
	\1 We have added a new subsection (Sec. III.A) to explain basic RMP 
	concepts.
	\1 We have also improved the flow of information of Sec. III to explain our 
	motivations to use this framework.
	\1 Sec. IV adds more details about our design and testing choices.
\end{outline}

\begin{displayquote}
- Some theoretical aspects of the approach should be discussed (see, in
particular, Reviewer 4);
\end{displayquote}

\begin{outline}[enumerate]
	\1 We have added an explanation for the choice of the tunable parameters 
	used by the controller in Sec. IV.
	\1 In this paper our focus has been in the application of the 
	RMP-based controller to provide safe local navigation, and have provided 
	experimental results in simulation as well as in real-world scenarios. A 
	further theoretical analysis of the stability and convergence properties 
	was beyond the scope of this paper.
	\1 Nevertheless, we have added a new Section (III.A) to discuss the RMP 
	framework in a more general way, and we have pointed related literature 
	that has discussed some of its theoretical properties [4],[30].
\end{outline}

\begin{displayquote}
- The contributions of the work need to be clarified.
\end{displayquote}

\begin{outline}[enumerate]
	\1 The main contribution of this work is the design, implementation and 
	deployment of a reactive controller for navigation tasks.
	\1 In particular, we have shown how to design a RMP controller that 
	exploits local maps available in legged robots, which is able to generate 
	reactive behaviors on a CPU.
	\1 These contributions have been reviewed and stated in the introduction 
	section.
	\1 Further, we have provided further explanation in the related works 
	section by placing our work in the context of other related methods in the 
	literature to make the motivation, differences, and contributions more 
	evident.
\end{outline}

%\begin{displayquote}
%Therefore, I believe that this paper should be revised considering the
%Reviewers' comments.
%\end{displayquote}


%%%%%%%%%%%%%%%%%%%%%%%%%%%%%%%%%%%%%%%%%%%%%%%%%%%%%%%%%%%%%%%%%%%%%%%%%%%%%%%%

\pagebreak
\section*{Response to Reviewer \#2 (ID 164917)}

\begin{displayquote}
	The paper presents an RMP-based approach to ensuring safe local
	movement for teach and repeat robots during their repeat traverse. The
	approach is (theoretically) generalized to any type of robot, rather
	than being specific to legged robots, and relies on some fairly strong
	(find to be out-of-scope for this paper) assumptions around terrain
	maps and other information. The system is set up to be fairly lean
	compute-wide and to run in real time on CPU. Some experiments are
	presented in simulation and in a short abandoned mine traverse. The
	paper is generally quite well written and presented and nicely concise
	in presenting the core theory.

	It's a little unclear throughout the paper what is simply standard best
	practice (e.g. in painting) and what is a new contribution specific to
	this paper. Presenting the key steps like in painting is fine for
	self-containedness of the paper of course.
\end{displayquote}

We thanks for the accurate summary and kind words.

\begin{outline}[enumerate]
	\1 As mentioned earlier, we have emphasized the contributions of our 
	work in the introduction. 
	\1 Inpainting is important in our pipeline because the later stages 
	require a well-defined image, so we considered it relevant for 
	self-containedness as it was pointed out.
	\1 The choice of in-painting technique typically depends on the noise 
	characteristics of the sensor and properties of the surface (eg. 
	reflectivity). In this paper, we use an inpainting techinique based on the 
	Fast Marching method presented by Telea [31], which is available in OpenCV. 
	\1 We preferred an inpainting technique because it produced a reasonable 
	guess for the unknown space at fast speed. We tried more advanced methods 
	(Gaussian Process regression) but they were too slow.
\end{outline}


\begin{displayquote}
	The RMP approach nicely combines various conflicting goals and
	constraints, which is neat.
	
	Most of the experiments focus on illustrative and qualitative
	evaluation and some parameter examination of the effects of various
	system configurations - so internal analysis. There's no benchmarking
	comparison to complete other methods, which admittedly would be
	difficult in this type of application. However, it would still help to
	be a little more explicit about what the current standards of
	performance are, and why what is presented here represents a
	substantial step beyond current capability in some (can be narrow)
	manner (compute, robustness etc...)
	
	For example, while not exactly teach and repeat, a bunch of robots just
	impressively solved similar problems in the DARPA SubT challenge (which
	this work may be associated with). How is what is presented here better
	or adds contributions beyond all that work? A direct comparison and
	re-implementation of those competitor methods (many of which probably
	aren't available) is probably unreasonable here, but there is likely
	some middle ground that provides some better semi-quantitative
	experimental comparison (even in a somewhat meta manner) of what the
	method provides experimentally over those others. 
\end{displayquote}


\begin{outline}[enumerate]
	\1 The lack of comparison to baseline methods has been pointed out in all
	the reviews. To address this aspect of the paper, we have now added a new 
	section for analyzing the performance with respect to other approaches.
	
	\1 In order to perform the comparison, we designed a simulated 
	environment (Fig. 11) with tight corners and varying corridor widths. The 
	environment has been designed in a manner that the different controllers 
	can be evaluated in challenging situations where a locally safe control 
	becomes critical.
	
	\1 We compared our RMP based controller against (a) FALCO-based~[31] 
	local planner (which is DARPA SubT related), (b) Potential field 
	controller, and (c) GDF based controller. This was added to Sec. IV and the 
	video.
	
	\1 Our experiments showed a consistent behavior of our proposal compared to 
	our real experiments, and the advantages of using potential field-like 
	methods when performing omnidirectional motion without requiring more 
	expensive planners.
\end{outline}

\begin{displayquote}
	Further work - how much could this framework be adapted beyond the
	teach and repeat context to general environment exploration, and would
	it retain is claimed benefits? I can't see anything obvious making this
	particularly limited to repeat in T\&R - frontier-based exploration with
	a local elevation map from LIDAR etc... should be enough to work in
	other applications?
\end{displayquote}

\begin{outline}[enumerate]
	\1 Indeed, while the controller was developed inspired by VT\&R-specific 
	requirements, the controller presented in this work can be used out that 
	context as well, since it only requires (1) a state estimate in a fixed 
	frame, 
	(2) an elevation map, and (3) a low-level controller that receives twists 
	commands.
	\1 For example, one of the ways we use this controller in our lab is simply 
	as a safety layer while an operator joysticks the robot.
	This controller ensures that even if the human operator provides a bad 
	control input which would result in a collision, the safety controller 
	would ensure that such a command is not executed. This greatly improves the 
	usage of the platform in general.
\end{outline}

%%%%%%%%%%%%%%%%%%%%%%%%%%%%%%%%%%%%%%%%%%%%%%%%%%%%%%%%%%%%%%%%%%%%%%%%%%%%%%%%

\pagebreak
\section*{Response to Reviewer \#3 (ID 164919)}

\begin{displayquote}
	The main contribution of the presented paper is an obstacle avoidance
	controller which uses Riemannian Motion Policies. Other contributions
	are the integration of said controller into a visual teach and repeat
	system and its evaluation on the ANYmal C quadruped.

	The authors also claim the use of vector field representations for
	navigation as a contribution, but such representations have been used
	for navigation before, as the authors correctly state in the Related
	Work section [17, 25, 26, 27]. Consequently, this should be removed
	from the list of contributions.
\end{displayquote}


\begin{outline}[enumerate]
	\1 Vector field representations have been effectively used for navigation 
	tasks in the past, but mainly to replace planners and in offline settings.
	\1 In this work, our contribution has been in efficiently computing 
	these vector fields online, and combining them in a RMP based controller 
	which runs at 10 Hz on a CPU.
	\1 To avoid the impression that the vector field representations themselves 
	are a contribution of this work, we have re-worded the second contribution 
	sentence in the Introduction section.
\end{outline}

\begin{displayquote}
	The Related Work section mentions important and relevant work and
	generally places the presented paper well in the existing literature.
	However, when distinguishing this work from Bista et al. [30] the
	authors claim that taking into account 3D geometry is the main unique
	feature. This, however, is not the case. The problem formulation of the
	Riemannian Motion Policy (RMP) is in SE2 and navigation cues are taken
	from 2-dimensional grid maps. The presented paper only uses 3D geometry
	(which actually isn’t really 3D, because it’s a 2.5D height map) to
	compute traversability values, which are thresholded to obtain an
	obstacle map. This is practically the same as [30], where they use
	depth images to compute an occupancy map around the robot. Having said
	this, this does not invalidate the contributions mentioned above, but
	the authors should more accurately and more clearly distinguish the
	presented RMP formulation from the heuristic on rotational velocity
	presented by Bista et al. [30]. 
\end{displayquote}

Thanks for this observation. We have done a few modifications to address this:
\begin{outline}[enumerate]
	\1 We have reworded that sentence to make the differences with Bista et al. 
	more explicit.
	\1 In particular, at the localization level, our representations and 
	localization strategies are different. We use topo-metric, with explicit 
	pose estimation, while they use purely topological and a visual-servoing 
	strategy.
	\1 In terms of the controller, as mentioned by the reviewer, Bista et al. 
	uses a heuristic to combine visual servoing's output with the A$\star$ 
	plan, while we exploit the RMP framework to combine different desired 
	behaviors in a principled way.
\end{outline}


\begin{displayquote}
	The paper is generally well written and well organized. The
	introduction motivates the work well and clearly states the scope and
	purpose of the presented method. Figures are clear and of high quality.
	Some minor grammatical errors should be fixed (e.g. missing words,
	“quadruped robot” $\rightarrow$ “quadrupedal robot”).
	The attached video illustrates the method very well (also, cute carrot
	:-) ).
\end{displayquote}

Thanks for this feedback, it is very much appreciated. 
\begin{outline}[enumerate]
	\1 We have corrected the typos mentioned in the reviews.
	\1 We also did another round of proof-reading and corrections.
\end{outline}
 

\begin{displayquote}
	The paper is generally technically sound, however, a few issues and
	questions do arise:
	
	The authors claim “In an environment without obstacles, the GDF would
	suffice to generate a motion policy to guide the robot towards the
	goal. However, just relying on the GDF is insufficient and other
	priorities such as moving in a certain direction or avoiding obstacles
	need to be considered.”. However, the computation of the GDF already
	takes traversability information, and therefore obstacles, into
	account. The use of additional SDF terms for obstacles avoidance
	instead of just inflating the obstacle map, should be motivated better. 
\end{displayquote}

We agree that this sentence was confusing and have updated the explanation in Sec. III.
We address the use of SDF terms for obstacle avoidance in the following way:
\begin{outline}[enumerate]
	\1 We agree that the GDF implicitly takes into account the obstacles. 
	However, it is only defined for the free space but it has discontinuities 
	near the obstacles, which results in a gradient that is not smooth.	
	\1 Additionally, its gradient does not necessarily guide the robot 
	to safer areas, but just tries to \emph{push} the robot toward the goal.
	\1 Inflating and smoothing the map is not straightforward for elongated 
	robots because the robot is longer along one axis, so there can still be 
	feasible solutions if the robot rotates in place. 
	\1 By formulating obstacle avoidance as a result of an independent field (here by SDF term), 
	we have more flexibility to control the desired behavior.
	\1 Lastly, in the new baselines experiment we have shown the resulting 
	trajectory when using a GDF-only approach.
\end{outline}

\begin{displayquote}
	In the first experimental run, the carrot point is set 0.7m from the
	robot, however the $d_c$ cutoff is 1m. Therefore, it seems like during
	this run the GDF is never active, since GDF is also disabled close to
	the goal. Is it required then, since the path tracking performance is
	similar to later runs with 1.5m carrot point distance? You mention
	speed as the main difference, was there any other qualitative
	difference in behavior?
\end{displayquote}

This is an important observation that we would like to clarify.
\begin{outline}[enumerate]
	\1 The distances $d_c$ are not \emph{hard} cuttofs but soft, since they are 
	the result of a logistic weighting (Tab. I, footnote).
	\1 Consequently, moving the carrot closer does not disable the GDF 
	entirely, but it gives it a lower weight compared to the Obstacle-free Goal 
	Reaching term.
	\1 As a result, the robot tries to minimize the Euclidean distance to the 
	goal instead of the geodesic one, while still having some geodesic terms 
	pulling the robot around the obstacle.
	\1 Without the contribution of the GDF term, the controller would locally 
	behave as a potential field controller and would get stuck in the obstacle 
	if the forces cancel each other.
	\1 However, since there is still a GDF term, we observed that when the robot 
	gets close to an obstacle, the GDF term still conservatively ``pulls'' 
	the robot around the obstacle. This also explains the lower average speed 
	for the \emph{Forward1} experiment in Sec. IV C.
\end{outline}

\begin{displayquote}
	Why does the goal-reaching term have a translational component?
	Wouldn’t the GDF term behave equally to the goal-reaching term in an
	obstacle-free environment and perform the same job? Why do you need to
	switch then, instead of just enabling a rotational goal-reaching
	component? The simulation experiment section states that the only issue
	with GDF is the lack of rotational component.
\end{displayquote}

\begin{outline}[enumerate]
	\1 This observation is correct, and we could effectively change the 
	\emph{Obstacle-free Goal Reaching} term for an orientation-only RMP to 
	achieve a similar behavior if we want to simplify the controller.
	\1 However, switching the RMPs via metrics is still required, since it is 
	desired to walk forward most of the time and only correct orientation to 
	match the goal when the robot is close.
\end{outline}

\begin{displayquote}
	Table 1 states that the heading correction is only active close to the
	goal. Shouldn’t it be the opposite? Far away from the goal you want to
	align with the motion direction, close to the goal you want to have the
	goal direction.
\end{displayquote}

Thanks for spotting this, it was a typo. It has been corrected in the new 
version.

\begin{displayquote}
	What are the effects of damping and regularization? How does the method
	perform without them? How do you choose the weighting constants $k_b$ and 	
	s?
\end{displayquote}

\begin{outline}[enumerate]
	\1 The damping term is usually used in RMP literature [4],[30] to control 
	the 
	acceleration profile and avoid overshooting. In particular, the quotient 
	$k_{\text{goal}} / k_{\text{damping}}$ determines the maximum speed 
	assuming no other accelerations are considered.
	\1 The regularization term puts a check on the difference in acceleration 
	at time $t$ and $t-1$, resulting in a smoother control without large jerks, 
	as well as avoiding ill-conditions in the optimization during experiments 
	with different RMPs configurations.
	\1 We obtain the weighting constant $k_b$ as well as other gains used in 
	the controller with a line search performed offline. We have added a 
	sentence in the paper to state this.
\end{outline}

\begin{displayquote}
	The paper claims safe navigation. How are the obstacle avoidance RMPs
	guaranteed to be of higher importance than things like goal-reaching
	and regularization.
\end{displayquote}

\begin{outline}[enumerate]
	\1 While we do not provide any formal guarantees for the controller, we 
	have tested the approach in challenging environments, both in simulation 
	as well as in real-world tests.
	\1 As we shown in the new Experiment 3, we leverage the effect of the 
	state-dependent Riemannian metrics with the reactivity of the controller to 
	dynamically change the desired behavior.
	\1 As we mentioned before, the parameters were chosen via line search in 
	simulation, and validated when we deployed the system.
	\1 Still, our implementation allows to easily change parameters on the fly, 
	which allows easy re-tuning of the parameters when testing the system in 
	new environments.
	%and have not observed any situations in which
	%the robot collided with obstacle as a result of 
	%thgoal-reaching/regularization 
	%forces dominating. 
	%\1 In particular, 
	%\1 In addition to the Riemannian metric which ensures the forces are 
	%scaled 
	%appropriately based on the distance 
	%to the obstacle, we also choose the gains for each RMP via a line search, 
	%which provides desired behavior for the robot. 
\end{outline}

\begin{displayquote}
	How does this approach compare to other, similar methods which use the
	same environment representation (SDF, GDF) but without RMPs. The
	experimental section demonstrates that this approach works, but it is
	hard to judge whether other methods could perform similarly.
\end{displayquote}

We have added a new experiment (Sec. IV.D) to compare against a potential field 
controller and a GDF-only controller, which is also shown in the supplementary 
video.

\begin{displayquote}
	It is difficult to judge how good or bad the path tracking error is.
	First, because there is a minimal tracking error due to the presence of
	obstacles. A “minimal error” or similar could better help judge the
	performance of your work. Second, comparing to other methods would be
	helpful to judge whether your method is not only safer, but also has
	better tracking accuracy.
\end{displayquote}

We have added a few sentences to clarify this:

\begin{outline}[enumerate]
	\1 Sec. III.D briefly discusses the expected performance of the VT\&R 
	system as reported in our previous work [5].
	\1 Experiment 1 briefly discusses the PTE obtained in the experiment 
	compared to the original VT\&R.
	\1 In general the system's tracking accuracy was within the expected bounds 
	but also being able to stay in safe areas without collisions and properly 
	recovering after losing track of the reference path.
\end{outline}
%%%%%%%%%%%%%%%%%%%%%%%%%%%%%%%%%%%%%%%%%%%%%%%%%%%%%%%%%%%%%%%%%%%%%%%%%%%%%%%%

\pagebreak
\section*{Response to Reviewer \#4 (ID 164921)}

\begin{displayquote}
	This paper develops a reactive control scheme for VT\&R missions. The
	controller acts as a safety filter on the reference trajectories
	computed by the VT\&R algorithm, relying on the information provided by
	two vector-field representations of the local environment, which are
	combined into a suitable RMP formulation. The controller is validated
	experimentally on a robot in two real-world environments. Overall, the
	paper is well-written. The method is interesting and the obtained
	results seem to be promising. However, some aspects of the presentation
	should be further clarified (see detailed comments below) and the paper
	lacks a comparison with other local guidance methods (e.g., local
	planners).
	
	Detailed comments:
	Fig.2: Please close the loop with the environment.
\end{displayquote}

Thanks for the valuable feedback. Regarding the comment, we agree that to 
properly illustrate the closed-loop system we should close the loop in the 
figure, but it was designed this way due to space limitations.

\begin{displayquote}
	Section III.B: The discussion here is too general and fails to
	provide a clear link with the control system architecture in Fig. 2 and
	the RMPs in Table I. Please expand this section a bit. For consistency,
	consider replacing the notation $v_B$ and $a_B$ with $\dot{x}_B$ and
	$\ddot{x}_B$. Perhaps, Section III.C can be moved before this section (it
	describes an input to the controller).
\end{displayquote}

Thanks for raising this concern. We have done a few modifications to connect 
the ideas better and improve the presentation:
\begin{outline}[enumerate]
	\1 We reordered the explanation of Figure 2 at the start of Sec. III.
	\1 We added a new section (Sec. III.A) to explain RMPs and the 
	motivation to use them as suggested. 
	\1 The explanation in Sec. III. C (formerly III.B) now is focused in the 
	specific details to implement the RMP framework, connecting it with Fig. 
	2, Tab. I and the rest of the pipeline.
	\1 Regarding the notation, in the general formulation (Eq. 1) we 
	use $\mathbf{x}$, $\mathbf{\dot{x}}$, and $\mathbf{\ddot{x}}$ to express a 
	generic state, velocity and acceleration. In the formulation we implemented 
	(Eq. 2) we changed the 
	notation to $\notationMatrixFrame{T}{}{\notationFrame{IB}}$ for the state, 
	$\notationVectorFrame{v}{B}{}$ for the velocity, and 
	$\notationVectorFrame{a}{B}{}$ for the acceleration to make explicit that 
	the state lies in $\SEtwo$ while the others are in $\Rn{3}$, avoiding 
	notation like $\notationMatrixFrame{\dot{T}}{}{\notationFrame{IB}}$ which 
	could be misleading.
	
\end{outline}

\begin{displayquote}
	The proposed architecture is validated empirically in Section IV, but
	many theoretical (and practically relevant) issues remain largely
	unaddressed. Does the control system enforce closed-loop stability
	(under reasonable assumptions)? Are collision avoidance and convergence
	to the goal guaranteed? How are the parameters $d_{carrot}$ and $d_c$ tuned?
	Please add some comments on these points.
\end{displayquote}


\begin{outline}[enumerate]
	\1 We have added an explanation for the choice of the tunable parameters 
	used by the controller in Sec. IV.
	\1 In this paper our focus has been in the application of the 
	RMP-based controller to provide safe local navigation, and have provided 
	experimental results in simulation as well as in real-world scenarios. A 
	further theoretical analysis of the stability and convergence properties 
	was beyond the scope of this paper.
	\1 Nevertheless, we have added a new Section (III.A) to discuss the RMP 
	framework in a more general way, and we have pointed related literature 
	that has discussed  some of its theoretical properties [4],[30].
\end{outline}

\begin{displayquote}
	The impact of the VT\&R signal quality on the overall control
	performance should be further investigated.
\end{displayquote}

\begin{outline}[enumerate]
	\1 The VT\&R system builds upon our previous work (Ref. [5]), which was 
	mainly evaluated in obstacle-free environments. We showed 
	that the path tracking system achieved less than \SI{10}{\centi\meter} 
	error w.r.t the reference path on average.
	\1 We have added a sentence stating this in the paper clarify the expected 
	accuracy of the VT\&R signal.
	\1 We also briefly discuss this in Section IV, Experiment 1.
\end{outline}

\begin{displayquote}
	Page 2, end of first column: please add a reference to Geodesic
	Distance Fields (GDF).
\end{displayquote}

We added a reference as requested.

\begin{displayquote}
	Section III.A, point 4: "a global minima", does this refer to a
	unique global minimum?
\end{displayquote}

Thanks for spotting this typo, it has been fixed.

\begin{displayquote}
	Table I, fourth column (obstacle avoidance), fourth row (when
	enabled?): from my understanding, this entry should be changed to
	"always".
\end{displayquote}

\begin{outline}[enumerate]
	\1 The table is correct by stating that the \emph{Obstacle Avoidance} RMPs 
	are enabled only when the robot is closer than a certain distance.
	\1 If they were always enabled, the resulting controller would behave 
	similarly to a potential field controller, and could consequently generate 
	local minima as a result of obstacle forces pushing the robot away even 
	when it is far, which is problematic when entering narrow passages.
	\1 In contrast, by relying on a distance-enabled metric we can control the 
	behavior of this repulsion forces in a more flexible way.
	\1 This can be observed in the attached video, which includes new 
	experiments and comparisons with potential field methods.
\end{outline}

%%%%%%%%%%%%%%%%%%%%%%%%%%%%%%%%%%%%%%%%%%%%%%%%%%%%%%%%%%%%%%%%%%%%%%%%%%%%%%%%

\pagebreak
\section*{Response to Reviewer \#8 (ID 167635)}

\begin{displayquote}
	The presents an efficient reactive collision avoidance system for safe
	navigation for quadrupedal robots. More specifically, the system tries
	to find a proper twist command from the given elevation map. The paper
	approaches this by combining multiple potential fields using a
	Riemannian Motion Policy framework, proposed by Ratliff et al. Then the
	proposed algorithm is combined with the quadrupedal robot, Anymal-C,
	and tested on indoor and cave navigation scenarios. The proposed
	algorithm successfully manages to avoid collisions while closely
	following the given path.

	Surprisingly, I am not too positive about the paper, even with the
	amazing hardware experiments, due to two main concerns. The first
	concern is its novelty: the proposed algorithm framework with RMPs is
	based on the existing work without any modification and never compared
	against other baseline algorithms. Second, the paper presents many
	concepts, such as VT\&R or RMPs, without any proper introductions, which
	makes readers hard to follow. Particularly, RMPs seems to be a fairly
	new concept introduced in 2018, which should not be expected to know.
	Therefore, I would like to recommend another round of revision to this
	paper, although it has many impressive results.
\end{displayquote}

Thanks for the feedback, we hope that the new modifications address the 
concerns:
\begin{outline}[enumerate]
  \1 While the RMP approach has been used for several applications including 
  	 navigation and control tasks [4],[30]
	 our contribution in this paper has been in applying the framework to design an effective controller which leverages local
	 perception to provide a safe navigation capability, as well as deploy and  
	 evaluate its performance on a quadrupedal robot in real-world experiments. 
  \1 In order to give the reader unfamiliar with RMP, we have improved our 
  	 motivation to use RMPs in Sec. II., and attached one additional reference 
  	 [30].
  \1 Additionally, we added a new section (III.A) to briefly give a general 
  	 overview of RMPs.
  \1 Now the original explanation we had in Sec. III.C (formerly III.B) is 
  	 focused on the specific details of our RMP implementation. 
  \1 We hope that these new modifications improve the readibility of the paper 
 	 in spite of the strict page-limit that difficults a more detailed 
 	 introduction to RMPs.  
\end{outline}


\begin{displayquote}
	I believe the idea of using potential fields to handle collision
	avoidance is not a new idea, as also suggested in the paper. The paper
	proposes a bit fancier formulation using RMP, but the core idea seems
	the same (at least not well highlighted in the paper). To make the
	matter worse, the algorithm is not compared with any existing
	algorithms - I believe there will be a lot of collision avoidance
	algorithms, even if we limit the scope to the potential fields based
	approaches. I would suggest the authors take more in-depth analysis to
	the given framework, such as ablation study.
\end{displayquote}

We have added a new experiment (Sec. IV, Experiment 3) to address this concern:

\begin{outline}[enumerate]
	\1 Comparisons are performed against a potential field method, a GDF-only 
	approach, and a local planner that determines a path via collision checking 
	and tracks it using a P controller.
	\1 Due to page limits, we have added Fig. 11 to show the main findings of 
	this experiment, and the full runs are shown in the supplementary video.
	\1 We also added a discussion about the behaviors we observed for each 
	method.
\end{outline}

\begin{displayquote}
	Another main concern is about the exposition. The field of robotics is
	very broad. But the authors used a few terms without explanation, such
	as VT\&R and RMP. Particularly, I am concerned about the RMP which is a
	relatively new idea introduced in 2018. I believe it is much safer to
	assume that readers are not familiar with this. The introduction does
	not convey the high level idea, and the explanation in Section III.B is
	quite shallow. Fig. 5 is hard to interpret and not sure about its
	message. Therefore, I believe the manuscript requires a major revision
	to fix the current exposition issue.
\end{displayquote}

As we have mentioned in the second response, we have reorganized the exposition 
of the method in Sec. III. We hope that this improves the readability of the 
paper and as well as the understanding of the RMP framework considering the 
space limitations.

% \pagebreak
% 
% \section*{Changes to the Video}
% 
% Below are the amendments to the video:
% \begin{itemize}
% \end{itemize}

\pagebreak

\includepdf[pages=-]{diff.pdf}

\end{document}
